\chapter{Fourier Series and Harmonic Analysis}
\label{chap:fourier-series}

\section{Introduction to Periodic Functions}
Periodic phenomena occur ubiquitously throughout electrical signal processing, vibration mechanics, and heat conduction. A function $f(x)$ is periodic with period $T > 0$ if:
\begin{equation}
f(x + T) = f(x), \quad \forall x \in \mathbb{R}
\end{equation}

\begin{theorem}[Orthogonality of Trigonometric Functions]
On the interval $[-\pi, \pi]$:
\begin{align}
\int_{-\pi}^{\pi} \cos(nx) \cos(mx) \, dx &= \begin{cases} 0, & n \neq m \\ \pi, & n = m \neq 0 \\ 2\pi, & n = m = 0 \end{cases} \\
\int_{-\pi}^{\pi} \sin(nx) \sin(mx) \, dx &= \begin{cases} 0, & n \neq m \\ \pi, & n = m \neq 0 \end{cases} \\
\int_{-\pi}^{\pi} \cos(nx) \sin(mx) \, dx &= 0, \quad \forall n, m
\end{align}
\end{theorem}

\section{Fourier Series and Euler's Formulae}
\begin{definition}[{Fourier Series Expansion in $[-\pi, \pi]$}]
Let $f(x)$ be defined on $[-\pi, \pi]$ with period $2\pi$. Its Fourier series is:
\begin{equation}
f(x) = \frac{a_0}{2} + \sum_{n=1}^{\infty} \left[ a_n \cos(nx) + b_n \sin(nx) \right]
\end{equation}
where the Euler-Fourier coefficients are given by:
\begin{align}
a_0 &= \frac{1}{\pi} \int_{-\pi}^{\pi} f(x) \, dx \\
a_n &= \frac{1}{\pi} \int_{-\pi}^{\pi} f(x) \cos(nx) \, dx, \quad n = 1, 2, 3, \dots \\
b_n &= \frac{1}{\pi} \int_{-\pi}^{\pi} f(x) \sin(nx) \, dx, \quad n = 1, 2, 3, \dots
\end{align}
\end{definition}

\begin{theorem}[Dirichlet's Conditions for Convergence]
Any function $f(x)$ can be expanded in a convergent Fourier series if, within the interval of definition:
\begin{enumerate}
    \item $f(x)$ is periodic, single-valued, and bounded.
    \item $f(x)$ has a finite number of finite discontinuities in any one period.
    \item $f(x)$ has a finite number of local maxima and minima in any one period.
\end{enumerate}
At a point of finite jump discontinuity $x = x_0$, the Fourier series converges to the arithmetic mean:
\begin{equation}
S(x_0) = \frac{f(x_0^+) + f(x_0^-)}{2}
\end{equation}
\end{theorem}

\section{Functions of Arbitrary Period 2L (Change of Interval)}
For a function defined on $[-L, L]$ of period $2L$:
\begin{equation}
f(x) = \frac{a_0}{2} + \sum_{n=1}^{\infty} \left[ a_n \cos\left(\frac{n\pi x}{L}\right) + b_n \sin\left(\frac{n\pi x}{L}\right) \right]
\end{equation}
where:
\begin{align}
a_0 &= \frac{1}{L} \int_{-L}^{L} f(x) \, dx \\
a_n &= \frac{1}{L} \int_{-L}^{L} f(x) \cos\left(\frac{n\pi x}{L}\right) \, dx \\
b_n &= \frac{1}{L} \int_{-L}^{L} f(x) \sin\left(\frac{n\pi x}{L}\right) \, dx
\end{align}

\section{Even and Odd Functions}
Symmetry drastically simplifies Fourier coefficient computation:
\begin{enumerate}
    \item \textbf{Even Functions ($f(-x) = f(x)$)}: The graph is symmetric about the $y$-axis.
    \begin{equation}
    b_n = 0, \quad a_0 = \frac{2}{L}\int_0^L f(x)\,dx, \quad a_n = \frac{2}{L}\int_0^L f(x)\cos\left(\frac{n\pi x}{L}\right)\,dx
    \end{equation}
    The series contains only cosine terms.
    \item \textbf{Odd Functions ($f(-x) = -f(x)$)}: The graph is symmetric about the origin.
    \begin{equation}
    a_0 = 0, \quad a_n = 0, \quad b_n = \frac{2}{L}\int_0^L f(x)\sin\left(\frac{n\pi x}{L}\right)\,dx
    \end{equation}
    The series contains only sine terms.
\end{enumerate}

\section{Half-Range Expansions}
When $f(x)$ is defined only on the half-interval $[0, L]$:
\begin{itemize}
    \item \textbf{Half-Range Cosine Series}: By extending $f(x)$ as an even function onto $[-L, 0]$:
    \begin{equation}
    f(x) = \frac{a_0}{2} + \sum_{n=1}^{\infty} a_n \cos\left(\frac{n\pi x}{L}\right), \quad a_n = \frac{2}{L}\int_0^L f(x)\cos\left(\frac{n\pi x}{L}\right)\,dx
    \end{equation}
    \item \textbf{Half-Range Sine Series}: By extending $f(x)$ as an odd function onto $[-L, 0]$:
    \begin{equation}
    f(x) = \sum_{n=1}^{\infty} b_n \sin\left(\frac{n\pi x}{L}\right), \quad b_n = \frac{2}{L}\int_0^L f(x)\sin\left(\frac{n\pi x}{L}\right)\,dx
    \end{equation}
\end{itemize}

\begin{example}[Fourier Series of a Square Wave]
Find the Fourier series of $f(x) = \begin{cases} -k, & -\pi < x < 0 \\ k, & 0 < x < \pi \end{cases}$, with period $2\pi$.

\textbf{Solution:}
Notice that $f(-x) = -f(x)$, so $f(x)$ is an odd function.
Consequently, $a_0 = 0$ and $a_n = 0$ for all $n \ge 1$.
Now calculate $b_n$:
\begin{align*}
b_n &= \frac{2}{\pi}\int_0^{\pi} k \sin(nx) \, dx = \frac{2k}{\pi} \left[ -\frac{\cos(nx)}{n} \right]_0^{\pi} \\
    &= \frac{2k}{n\pi} (1 - \cos n\pi) = \frac{2k}{n\pi} (1 - (-1)^n)
\end{align*}
For even $n$, $b_n = 0$. For odd $n = 2m-1$, $b_n = \frac{4k}{(2m-1)\pi}$.
Thus:
\begin{equation*}
f(x) = \frac{4k}{\pi}\left( \sin x + \frac{\sin 3x}{3} + \frac{\sin 5x}{5} + \dots \right)
\end{equation*}
\end{example}
